\section{Perfect complexes and K-groups}
\label{section-perfect-K-group}

\noindent
We quickly show that the zeroth K-group of the derived category of
perfect complexes of a ring $R$ is the same as $K_0(R)$ defined in
Algebra, Section \ref{algebra-section-K-groups}.

\begin{lemma}
\label{lemma-perfect-to-K-group}
Let $R$ be a ring. There is a map
$$
c : \text{perfect complexes over }R \longrightarrow K_0(R)
$$
with the following properties
\begin{enumerate}
\item $c(K[n]) = (-1)^nc(K)$ for a perfect complex $K$,
\item if $K \to L \to M \to K[1]$ is a distinguished triangle of
perfect complexes, then $c(L) = c(K) + c(M)$,
\item if $K$ is represented by a finite complex $M^\bullet$
consisting of finite projective modules, then
$c(K) = \sum (-1)^i[M_i]$.
\end{enumerate}
\end{lemma}

\begin{proof}
Let $K$ be a perfect object of $D(R)$. By definition we can represent
$K$ by a finite complex $M^\bullet$ of finite projective $R$-modules.
We define $c$ by setting
$$
c(K) = \sum (-1)^n[M^n]
$$
in $K_0(R)$. Of course we have to show that this is well defined,
but once it is well defined, then (1) and (3) are immediate.
For the moment we view the map $c$ as defined on complexes of
finite projective $R$-modules.

\medskip\noindent
Suppose that $L^\bullet \to M^\bullet$ is a surjective map
of finite complexes of finite projective $R$-modules.
Let $K^\bullet$ be the kernel. Then we obtain short exact
sequences of $R$-modules
$$
0 \to K^n \to L^n \to M^n \to 0
$$
which are split because $M^n$ is projective. Hence $K^\bullet$
is also a finite complex of finite projective $R$-modules and
$c(L^\bullet) = c(K^\bullet) + c(M^\bullet)$ in $K_0(R)$.

\medskip\noindent
Suppose given finite complex $M^\bullet$ of finite projective $R$-modules
which is acyclic. Say $M^n = 0$ for $n \not \in [a, b]$. Then we
can break $M^\bullet$ into short exact sequences
$$
\begin{matrix}
0 \to M^a \to M^{a + 1} \to N^{a + 1} \to 0, \\
0 \to N^{a + 1} \to M^{a + 2} \to N^{a + 3} \to 0, \\
\ldots \\
0 \to N^{b - 3} \to M^{b - 2} \to N^{b - 2} \to 0, \\
0 \to N^{b - 2} \to M^{b - 1} \to M^b \to 0
\end{matrix}
$$
Arguing by descending induction we see that $N^{b - 2}, \ldots, N^{a + 1}$
are finite projective $R$-modules, the sequences are split exact, and
$$
c(M^\bullet) = \sum (-1)[M^n] = \sum (-1)^n([N^{n - 1}] + [N^n]) = 0
$$
Thus our construction gives zero on acyclic complexes.

\medskip\noindent
It follows formally from the results of the preceding two
paragraphs that $c$ is well defined and satisfies (2). Namely,
suppose the finite complexes $M^\bullet$ and $L^\bullet$ of
finite projective $R$-modules represent the same object
of $D(R)$. Then we can represent the isomorphism by a map
$f : M^\bullet \to L^\bullet$ of complexes, see
Derived Categories, Lemma \ref{derived-lemma-morphisms-from-projective-complex}.
We obtain a short exact sequence of complexes
$$
0 \to L^\bullet \to C(f)^\bullet \to K^\bullet[1] \to 0
$$
see Derived Categories, Definition \ref{derived-definition-cone}.
Since $f$ is a quasi-isomorphism, the cone $C(f)^\bullet$ is
acyclic (this follows for example from the discussion in
Derived Categories, Section \ref{derived-section-canonical-delta-functor}).
Hence
$$
0 = c(C(f)^\bullet) = c(L^\bullet) + c(K^\bullet[1]) =
c(L^\bullet) - c(K^\bullet)
$$
as desired. We omit the proof of (2) which is similar.
\end{proof}

\noindent
The following lemma shows that $K_0(R)$ is equal to
$K_0(D_{perf}(R))$.

\begin{lemma}
\label{lemma-perfect-to-K-group-universal}
Let $R$ be a ring. Let $D_{perf}(R)$ be the derived category of
perfect objects, see Lemma \ref{lemma-perfect-ring-classical-generator}.
The map $c$ of Lemma \ref{lemma-perfect-to-K-group} gives an isomorphism
$K_0(D_{perf}(R)) = K_0(R)$.
\end{lemma}

\begin{proof}
It follows from the definition of $K_0(D_{perf}(R))$
(Derived Categories, Definition \ref{derived-definition-K-zero})
that $c$ induces a homomorphism $K_0(D_{perf}(R)) \to K_0(R)$.

\medskip\noindent
Given a finite projective module $M$ over $R$ let us denote
$M[0]$ the perfect complex over $R$ which has $M$ sitting in degree $0$
and zero in other degrees. Given a short exact sequence
$0 \to M \to M' \to M'' \to 0$ of finite projective modules
we obtain a distinguished triangle $M[0] \to M'[0] \to M''[0] \to M[1]$, see
Derived Categories, Section \ref{derived-section-canonical-delta-functor}.
This shows that we obtain a map $K_0(R) \to K_0(D_{perf}(R))$
by sending $[M]$ to $[M[0]]$ with apologies for the horrendous notation.

\medskip\noindent
It is clear that $K_0(R) \to K_0(D_{perf}(R)) \to K_0(R)$ is the identity.
On the other hand, if $M^\bullet$ is a bounded complex of finite projective
$R$-modules, then the the existence of the distinguished triangles
of ``stupid truncations''
(see Homology, Section \ref{homology-section-truncations})
$$
\sigma_{\geq n}M^\bullet \to \sigma_{\geq n - 1}M^\bullet \to
M^{n - 1}[-n + 1] \to (\sigma_{\geq n}M^\bullet)[1]
$$
and induction show that
$$
[M^\bullet] = \sum (-1)^i[M^i[0]]
$$
in $K_0(D_{perf}(R))$ (with again apologies for the notation).
Hence the map $K_0(R) \to K_0(D_{perf}(R))$ is surjective which
finishes the proof.
\end{proof}







